% Generated draft: requires mathematical and source review.
\documentclass[11pt]{article}
\usepackage[margin=25mm]{geometry}
\usepackage{fontspec}
\setmainfont{Latin Modern Roman}
\usepackage{amsmath,amssymb,mathtools,unicode-math}
\setmathfont{Latin Modern Math}
\usepackage{enumitem,needspace,xurl,hyperref}
\setlength{\parindent}{0pt}
\setlength{\parskip}{5pt}
\setlength{\emergencystretch}{3em}
\newcommand{\N}{\mathbb{N}}
\newcommand{\Z}{\mathbb{Z}}
\newcommand{\Q}{\mathbb{Q}}
\newcommand{\R}{\mathbb{R}}
\newcommand{\C}{\mathbb{C}}
\newcommand{\F}{\mathbb{F}}
\newcommand{\A}{\mathbb{A}}
\newcommand{\PP}{\mathbb{P}}
\newcommand{\E}{\mathbb{E}}
\newcommand{\eps}{\varepsilon}
\newcommand{\dd}{\,\mathrm{d}}
\begin{document}
\section*{UnsolvedMath ID: 2100606}
\textbf{Source problem code:} catalogue entry
\subsection{Definitions and mathematical statement}
In the usual regular-fibration setting, let an integrable Hamiltonian system on a symplectic $2n$-manifold have a regular-value region $B$ over which its fibers form a locally trivial bundle $\pi:M_B\to B$ of compact connected Lagrangian $n$-tori. The fiber homology groups define an integral local system $\Lambda$ of rank $n$, equivalently the period-lattice local system under the action-coordinate identification. Transition functions of the torus fibration can have both integral linear and translation parts.

The Duistermaat--Chern class, in the usual convention, is the translation or torsor obstruction
\[
c(\pi)\in H^2(B;\Lambda),
\]
where cohomology has the indicated local coefficients. It obstructs a global section of the torus bundle. The exact convention for the period lattice and the class should be confirmed against the cited source. This class must be distinguished from monodromy, which records the linear holonomy in $\operatorname{GL}(n,\Z)$, and from additional obstructions to a global Lagrangian section or global action--angle coordinates. Nontrivial monodromy alone is not a computation of nonzero $c(\pi)$.

The source request, attributed to S. Vũ Ngọc, is to find natural examples from classical mechanics for which this Duistermaat--Chern class is nonzero. A satisfactory specified example must identify the mechanical system, its regular torus-fibration region $B$, the coefficient local system, and a calculation establishing
\[
c(\pi)\ne0.
\]
The source does not define ``natural'', restrict the mechanical phase spaces or Hamiltonians, specify dimension, or say whether the class is meant only on the regular region or in an extension across singular fibers. Those choices remain open rather than being replaced by an arbitrary abstract torus bundle.

The supplied partial-progress summary claims positive examples in the spherical pendulum and semitoric systems, and suggests that only a systematic list remains. It gives no selected fibration, base, local-coefficient computation, or reference supporting that claim. Since these names can involve monodromy and other singular-fibration invariants, the claim does not by itself verify the requested nonzero Duistermaat--Chern class. No positive example or completed existence resolution is asserted by this draft.
\subsection{Short English statement}
Find concrete classical mechanical integrable systems whose regular torus fibrations have a nonzero Duistermaat--Chern obstruction. The class and the chosen base must be computed explicitly; monodromy alone is a different invariant. The supplied summary names the spherical pendulum and semitoric systems but gives no Chern-class calculation, so its claimed positive resolution needs review.
\subsection{Sources}
\begin{itemize}
\item Local source registry: \url{data/unsolvedmath/problems.json}, permanent UnsolvedMath ID $2100606$, source code catalogue entry, including the supplied PARTIAL-PROGRESS summary. Dataset attribution: Ulam AI and UnsolvedMath Contributors; source attribution and rights are retained.
\item \emph{Open Problems, Questions, and Challenges in Finite-Dimensional Integrable Systems} (2018), Problem 6.6, attributed to S. Vũ Ngọc: \url{https://arxiv.org/abs/1804.03737}. The metadata reports source-TeX extraction and July 29, 2026 access, with status and rights review flags; no source-paper author list is supplied in this record.
\item The supplied research summary mentions the spherical pendulum and semitoric systems but gives no authoritative example citation, precise regular fibration, or Duistermaat--Chern-class computation. No linked source was inspected for this draft.
\end{itemize}
\end{document}
